\subsection{Homogeneous closed loci in projective charts}
\label{subsec:closed-projective-charts}

The algebraic statements here hold over any commutative coefficient ring.
The analytic statements concern the normalized projective model over a trivially
valued field. Homogeneity is retained explicitly; no identification with arbitrary
scheme analytifications or gluing of coherent sheaves is assumed.
It cannot be dropped: $X_j-1$ dehomogenizes to zero, while a graded seminorm
assigns $X_j-1$ the maximum of the values of its two homogeneous components,
which is at least the value $1$ of its constant component.

\begin{lemma}[Localization of a homogeneous chart ideal]
\label{closed-charts:localized-ideal}
\lean{EnvelopeProperty.ProjBerk.coordinateChartIdeal,
EnvelopeProperty.ProjBerk.coordinateChartIdeal_localization_eq}
\leanok
\uses{charts:homogeneous-substitution,charts:localized-homogeneous-substitution}
Let $R$ be a commutative ring, $S=R[X_0,\ldots,X_n]$, and let $J\subseteq S$
be homogeneous for the standard total-degree grading. Put
$J_j=D_j(J)A_j\subseteq A_j=R[U_i:i\ne j]$.
In $S[X_j^{-1}]$, the ideals generated by the images of $J$ and of
$\beta_j(J_j)$ are equal.
\end{lemma}
\begin{proof}
For every homogeneous component $F_d$, the chart scaling identity gives
$\beta_j(D_j(F_d))X_j^d=F_d$ in the localization. Since $X_j$ is a unit,
each side generates the same principal ideal there.
If $F\in J$, homogeneity implies $F_d\in J$ for every $d$.
The identity $F=\sum_dF_d$ therefore places the image of $F$ in the ideal
generated by $\beta_j(J_j)$, proving one inclusion.
Conversely $D_j(F)=\sum_dD_j(F_d)$ for $F\in J$.
The scaling identity and invertibility of $X_j$ show that every
$\beta_j(D_j(F_d))$ belongs to the localized ideal generated by $J$.
Their sum does as well. These $D_j(F)$ generate $J_j$, proving the reverse inclusion.
\end{proof}

\begin{lemma}[Closed-locus evaluation in affine charts]
\label{closed-charts:vanishing}
\lean{EnvelopeProperty.ProjBerk.projectiveZeroLocus,
EnvelopeProperty.ProjBerk.mem_projectiveZeroLocus,
EnvelopeProperty.ProjBerk.isClosed_projectiveZeroLocus,
EnvelopeProperty.ProjBerk.coordinateChartMap_mem_zeroLocus_iff}
\leanok
\uses{closed-charts:localized-ideal,charts:homeomorphism,
def:analytic-zero-locus,lem:analytic-zero-locus-closed}
For a homogeneous ideal $J\subseteq k[X_0,\ldots,X_n]$, define its closed
projective-model locus by $Z_J=\{v:J\subseteq\operatorname{supp}(v)\}$.
This locus is closed. For every $v\in V_j$,
$v\in Z_J$ if and only if $c_j(v)$ vanishes on $J_j$.
\end{lemma}
\begin{proof}
The locus is an intersection of closed evaluation zero sets, hence is closed.
Extend $v$ to the localization at $X_j$. Vanishing of $v$ on $J$ is equivalent
to vanishing of the extension on the ideal generated by its image, because the
extension preserves the value of every polynomial and its support is an ideal.
The preceding lemma identifies this localized ideal with that generated by
$\beta_j(J_j)$. Vanishing on this ideal is exactly vanishing of the restriction
along $\beta_j$, namely $c_j(v)$, on $J_j$.
\end{proof}

\begin{theorem}[Quotient spectra give the closed projective charts]
\label{closed-charts:quotient-homeomorphism}
\lean{EnvelopeProperty.ProjBerk.closedChart,
EnvelopeProperty.ProjBerk.isClosed_closedChart,
EnvelopeProperty.ProjBerk.closedChartHomeomorph,
EnvelopeProperty.ProjBerk.closedChartHomeomorph_apply_coe,
EnvelopeProperty.ProjBerk.quotientChartHomeomorph,
EnvelopeProperty.ProjBerk.coordinateChartMap_quotientChartHomeomorph}
\leanok
\uses{closed-charts:vanishing,charts:homeomorphism,lem:analytic-affine-quotient}
For each homogeneous ideal $J$, the closed part of the $j$-th coordinate chart
$Z_J\cap V_j$ is homeomorphic to $(\Spec(A_j/J_j))^\an$.
This is the actual affine quotient spectrum, with its pointwise seminorm topology.
\end{theorem}
\begin{proof}
The chart homeomorphism restricts, by the vanishing equivalence, to a homeomorphism
between $Z_J\cap V_j$ and the affine seminorm zero locus of $J_j$ in
$(\Spec A_j)^\an$. The quotient-spectrum homeomorphism identifies this latter
zero locus with $(\Spec(A_j/J_j))^\an$. Composing it with the inverse restricted
chart homeomorphism gives the stated map. Its two inverse identities and continuity
follow from those of the component homeomorphisms.
\end{proof}
